%% =====================================================================
%% Supplementary material for
%% "Operator-Adaptive Partial Least Squares for Near-Infrared
%% Spectroscopy: Unified AOM/POP Selection, Covariance-Space SIMPLS,
%% and PLS-DA"
%% =====================================================================

\documentclass[a4paper,11pt]{article}

\usepackage[utf8]{inputenc}
\usepackage[T1]{fontenc}
\usepackage{amsmath,amssymb,amsthm}
\usepackage{mathtools}
\usepackage{bm}
\usepackage{booktabs}
\usepackage{graphicx}
\usepackage{xcolor}
\usepackage{algorithm}
\usepackage{algpseudocode}
\usepackage{hyperref}
\usepackage{url}
\usepackage{microtype}
\usepackage{enumitem}
\usepackage{caption}
\usepackage{subcaption}
\usepackage{tabularx}
\usepackage{longtable}
\usepackage{multirow}
\usepackage{array}
\usepackage{geometry}
\geometry{margin=2.4cm}

\hypersetup{
  colorlinks=true,
  linkcolor=blue!50!black,
  citecolor=blue!50!black,
  urlcolor=blue!50!black,
}

\newcommand{\R}{\mathbb{R}}
\newcommand{\softmax}{\mathrm{softmax}}
\newcommand{\diag}{\mathrm{diag}}
\newcommand{\identity}{\mathbf{I}}
\newcommand{\Xmat}{\mathbf{X}}
\newcommand{\Ymat}{\mathbf{Y}}
\newcommand{\Tmat}{\mathbf{T}}
\newcommand{\Pmat}{\mathbf{P}}
\newcommand{\Wmat}{\mathbf{W}}
\newcommand{\Qmat}{\mathbf{Q}}
\newcommand{\Bmat}{\mathbf{B}}
\newcommand{\Amat}{\mathbf{A}}
\newcommand{\Smat}{\mathbf{S}}
\newcommand{\Zmat}{\mathbf{Z}}
\newcommand{\xvec}{\mathbf{x}}
\newcommand{\yvec}{\mathbf{y}}
\newcommand{\zvec}{\mathbf{z}}
\newcommand{\rvec}{\mathbf{r}}
\newcommand{\wvec}{\mathbf{w}}
\newcommand{\tvec}{\mathbf{t}}
\newcommand{\pvec}{\mathbf{p}}
\newcommand{\qvec}{\mathbf{q}}
\newcommand{\bvec}{\mathbf{b}}
\newcommand{\AOM}{\textsc{AOM}}
\newcommand{\POP}{\textsc{POP}}
\newcommand{\PLS}{\textsc{PLS}}
\newcommand{\PLSDA}{\textsc{PLS-DA}}
\newcommand{\NIPALS}{\textsc{NIPALS}}
\newcommand{\SIMPLS}{\textsc{SIMPLS}}

\title{Supplementary material:\\
Operator-Adaptive Partial Least Squares for Near-Infrared Spectroscopy}
\author{Gr\'egory Beurier and the \texttt{nirs4all} contributors}
\date{}

\begin{document}
\maketitle

\renewcommand{\thesection}{S\arabic{section}}

%% =====================================================================
\section{Extended derivations}\label{app:math}
%% =====================================================================

\subsection{The covariance identity}\label{app:identity}
Let $\Xmat\in\R^{n\times p}$ and $\Ymat\in\R^{n\times q}$ be centered.
Let $\Amat\in\R^{p\times p}$ be a strict linear operator that acts on
row spectra by $\Xmat_b=\Xmat\Amat^\top$. Then
\[
(\Xmat_b)^\top \Ymat
= (\Xmat\Amat^\top)^\top \Ymat
= \Amat\,\Xmat^\top \Ymat.
\]
The proof is one line of matrix algebra and uses only the fact that
$\Amat$ does not depend on the rows of $\Xmat$. Two corollaries follow
immediately:
\begin{enumerate}[leftmargin=2em]
  \item Cross-covariance evaluation: $\Smat_b=\Amat\Smat$ where
        $\Smat=\Xmat^\top \Ymat$. This makes operator scoring
        $\mathcal{O}(p^2 q)$ at most, and $\mathcal{O}(pq)$ when
        $\Amat$ is banded (e.g.\ Savitzky-Golay or finite difference).
  \item Effective weight: if $\rvec_b$ is a weight vector in the
        transformed space (for example, the dominant left-singular
        vector of $\Smat_b$), then the same direction is represented
        in the original space by $\zvec=\Amat^\top \rvec_b$. This is
        what lets per-component policies map back to a single
        coordinate system for the final coefficient $\Bmat$.
\end{enumerate}

\subsection{Effective weights with operator change}\label{app:effective}
For the per-component selection policy, the operator $\Amat^{(a)}$ may
change with the component index $a$. Let $\rvec_a\in\R^p$ be the
dominant left-singular vector of $\Amat^{(a)}\Smat$. The original-space
effective weight at component $a$ is $\zvec_a=\Amat^{(a)\,\top}\rvec_a$.
The score is $\tvec_a=\Xmat\zvec_a$, the loading is
$\pvec_a=\Xmat^\top\tvec_a/\|\tvec_a\|^2$. After the rank-one
deflation $\Smat\leftarrow \Smat - \pvec_a(\pvec_a^\top \Smat)$, the
next operator $\Amat^{(a+1)}$ is selected on the deflated $\Smat$ and
the construction repeats.

A subtle but important point: orthogonality of subsequent components
must be enforced in the original space, not in any single transformed
space, because the transformed spaces change between components. This
is the reason we expose an \texttt{orthogonalization="original"} mode
and we set it as the automatic default for \POP and the soft policy.
For the \AOM-\textsc{global} policy, the operator does not change and
both modes coincide; we keep \texttt{transformed} as the default for
\textsc{global} for parity with the materialized reference.

\subsection{Soft mixture degeneracy}\label{app:soft}
Let $\lambda\in\Delta^{B-1}$, the simplex on the bank
$\mathcal{B}=\{\Amat_1,\dots,\Amat_B\}$. The covariance objective for
the first component, given a normalised weight $\rvec\in\R^p$, is
\[
\Phi(\lambda,\rvec)
\;=\;\Big\|\rvec^\top \Big(\sum_b \lambda_b \Amat_b\Big) \Smat\Big\|.
\]
For fixed $\rvec$, $\Phi(\cdot,\rvec)$ is a convex function on the
simplex (it is the absolute value of a linear function of $\lambda$),
hence its maximum is attained at a vertex. Conversely, for fixed
$\lambda$, the optimal $\rvec$ is the dominant left-singular vector of
$\sum_b \lambda_b \Amat_b\Smat$. Alternation between $\lambda$ and
$\rvec$ therefore drives $\lambda$ to a vertex, which is exactly
\AOM-\textsc{global}. Hence the soft policy collapses to hard global
selection under the covariance objective. Numerically richer behaviour
can in principle be obtained by adding a strictly concave entropy
penalty to $\Phi$, but we do not pursue that here, and we report
\textsc{soft} only as a research interface.

\subsection{Class-balanced coding equivalence}\label{app:plsda-coding}
Let $\Ymat^{(0)}$ be the standard one-hot indicator matrix
($Y^{(0)}_{ic}=1$ if $y_i=c$, $0$ otherwise) and let $\Ymat$ be the
class-balanced encoding of equation~(7) of the main text:
$Y_{ic}=1/\sqrt{\pi_c}$ if $y_i=c$, $0$ otherwise. Write
$\Ymat=\Ymat^{(0)}\,\diag(\pi^{-1/2})$ where $\pi^{-1/2}$ is the vector
of class-rescaling factors. Then the cross-covariance becomes
\[
\Xmat^\top \Ymat
\;=\; \Xmat^\top \Ymat^{(0)}\,\diag(\pi^{-1/2}),
\]
which uniformly scales the per-class columns of the cross-covariance.
This is exactly the rebalancing one obtains from the
\verb|class_weight="balanced"| option of \texttt{sklearn} logistic
regression on a per-row scale, but applied to the latent extraction
itself. Class-balanced coding does not change the linear span of the
latent space; it changes which directions are dominant when classes
are imbalanced.

\subsection{Logistic calibration as a fold-internal classifier}\label{app:calibration}
For each component count $K$, the latent-score matrix
$\Tmat=\Xmat\Zmat\in\R^{n\times K}$ is fixed by the chosen operator
sequence and the regression coefficient $\Bmat$. The logistic
calibrator fits $\Wmat_{\mathrm{lr}}\in\R^{K\times K}$ and
$\bvec_{\mathrm{lr}}\in\R^K$ by maximising the multinomial likelihood
on $\Tmat$ with L2 regularisation
$\frac{1}{2C}\|\Wmat_{\mathrm{lr}}\|_F^2$ where $C=1$ by default. Because
$\Tmat$ depends only on training-fold information, the calibrator is
trained on the same fold and never sees the test set. The fallback
$\softmax(\Tmat/\tau)$ shares the same scope: the temperature $\tau$
is fitted by minimising training-fold log loss and is held fixed at
inference time.

%% =====================================================================
\section{Extended tables}\label{app:tables}
%% =====================================================================

\subsection{Per-variant per-dataset smoke regression}
\begin{longtable}{lll}
\caption{Smoke benchmark: RMSEP per variant per dataset (3 regression
datasets, 11 variants).}\\
\toprule
Variant & Dataset & RMSEP \\
\midrule
\endfirsthead
\toprule
Variant & Dataset & RMSEP \\
\midrule
\endhead
\bottomrule
\endfoot
PLS-standard-numpy & ALPINE/ALPINE\_P\_291\_KS & 0.0766 \\
AOM-global-nipals-materialized-numpy & ALPINE/ALPINE\_P\_291\_KS & 0.0766 \\
AOM-global-nipals-adjoint-numpy & ALPINE/ALPINE\_P\_291\_KS & 0.0766 \\
POP-nipals-materialized-numpy & ALPINE/ALPINE\_P\_291\_KS & 0.0721 \\
POP-nipals-adjoint-numpy & ALPINE/ALPINE\_P\_291\_KS & 0.0721 \\
AOM-global-simpls-materialized-numpy & ALPINE/ALPINE\_P\_291\_KS & 0.0766 \\
AOM-global-simpls-covariance-numpy & ALPINE/ALPINE\_P\_291\_KS & 0.0766 \\
POP-simpls-materialized-numpy & ALPINE/ALPINE\_P\_291\_KS & 0.0804 \\
POP-simpls-covariance-numpy & ALPINE/ALPINE\_P\_291\_KS & 0.0804 \\
AOM-soft-simpls-covariance-numpy & ALPINE/ALPINE\_P\_291\_KS & 0.0767 \\
PLS-standard-numpy & AMYLOSE/Rice\_Amylose\_313\_YbasedSplit & 4.3290 \\
AOM-global-nipals-materialized-numpy & AMYLOSE/Rice\_Amylose\_313\_YbasedSplit & 4.3290 \\
AOM-global-nipals-adjoint-numpy & AMYLOSE/Rice\_Amylose\_313\_YbasedSplit & 4.3290 \\
POP-nipals-materialized-numpy & AMYLOSE/Rice\_Amylose\_313\_YbasedSplit & 3.9181 \\
POP-nipals-adjoint-numpy & AMYLOSE/Rice\_Amylose\_313\_YbasedSplit & 3.9181 \\
AOM-global-simpls-materialized-numpy & AMYLOSE/Rice\_Amylose\_313\_YbasedSplit & 4.3290 \\
AOM-global-simpls-covariance-numpy & AMYLOSE/Rice\_Amylose\_313\_YbasedSplit & 4.3290 \\
POP-simpls-materialized-numpy & AMYLOSE/Rice\_Amylose\_313\_YbasedSplit & 4.4407 \\
POP-simpls-covariance-numpy & AMYLOSE/Rice\_Amylose\_313\_YbasedSplit & 4.4407 \\
AOM-soft-simpls-covariance-numpy & AMYLOSE/Rice\_Amylose\_313\_YbasedSplit & 4.2924 \\
PLS-standard-numpy & BEER/Beer\_OriginalExtract\_60\_KS & 0.5713 \\
AOM-global-nipals-materialized-numpy & BEER/Beer\_OriginalExtract\_60\_KS & 0.5713 \\
AOM-global-nipals-adjoint-numpy & BEER/Beer\_OriginalExtract\_60\_KS & 0.5713 \\
POP-nipals-materialized-numpy & BEER/Beer\_OriginalExtract\_60\_KS & 0.4212 \\
POP-nipals-adjoint-numpy & BEER/Beer\_OriginalExtract\_60\_KS & 0.4212 \\
AOM-global-simpls-materialized-numpy & BEER/Beer\_OriginalExtract\_60\_KS & 0.5713 \\
AOM-global-simpls-covariance-numpy & BEER/Beer\_OriginalExtract\_60\_KS & 0.5713 \\
POP-simpls-materialized-numpy & BEER/Beer\_OriginalExtract\_60\_KS & 0.4207 \\
POP-simpls-covariance-numpy & BEER/Beer\_OriginalExtract\_60\_KS & 0.4207 \\
AOM-soft-simpls-covariance-numpy & BEER/Beer\_OriginalExtract\_60\_KS & 0.5199 \\
\end{longtable}

\subsection{Smoke classification, partial run}
\begin{table}[h]
\centering
\caption{Smoke classification on \texttt{ARABIDOPSIS\_CEFE/Genotype10\_250}. Other
classification datasets in the smoke run hit a numerical
non-convergence and are reported as exception rows in
\texttt{results\_classification.csv}.}
\begin{tabular}{lrrrr}
\toprule
Variant & Bal.\ Acc. & Macro-F1 & log-loss & ECE \\
\midrule
PLS-DA-standard & 0.247 & 0.224 & 2.199 & 0.047 \\
AOM-PLS-DA-global-nipals-adjoint & 0.247 & 0.224 & 2.199 & 0.047 \\
POP-PLS-DA-nipals-adjoint & 0.154 & 0.127 & 2.243 & 0.030 \\
AOM-PLS-DA-global-simpls-covariance & 0.274 & 0.247 & 2.265 & 0.169 \\
POP-PLS-DA-simpls-covariance & 0.208 & 0.166 & 2.263 & 0.118 \\
\bottomrule
\end{tabular}
\end{table}

%% =====================================================================
\section{Dataset cards}\label{app:datasets}
%% =====================================================================

\subsection{Regression: ALPINE}
ALPINE is a NIRS quantification cohort included in the TabPFN/NIRS
benchmark. Train/test split is Kennard-Stone~\cite{kennard1969computer}
(KS) on the predictor space. The dataset has $\sim$291 samples and
moderate spectral dimensionality. Reference RMSEPs (master CSV):
\PLS $\approx 0.062$, TabPFN-Raw $\approx 0.060$,
TabPFN-opt $\approx 0.043$.

\subsection{Regression: AMYLOSE (Rice amylose)}
Rice amylose calibration with a $y$-based split. $\sim$313 samples;
the response variable is a positive concentration. Reference RMSEPs:
\PLS $\approx 1.91$, TabPFN-Raw $\approx 4.51$,
TabPFN-opt $\approx 1.63$.

\subsection{Regression: BEER}
BEER calibration of original extract concentration with $\sim$60
samples and a Kennard-Stone split. Reference RMSEPs: \PLS $\approx 0.38$,
TabPFN-Raw $\approx 0.25$, TabPFN-opt $\approx 0.13$.

\subsection{Classification: ARABIDOPSIS CEFE genotype}
Multi-class genotype classification with $\sim$250 samples and 10
classes (Genotype10\_250). Used as a controlled, balanced toy cohort for
the classification protocol.

%% =====================================================================
\section{Operator parameter choices}\label{app:operators}
%% =====================================================================

\subsection{Whittaker smoother}
The Whittaker smoother~\cite{eilers2003perfect} solves
$(\identity+\lambda \mathbf{D}^\top \mathbf{D})\,\xvec=\xvec$ where
$\mathbf{D}$ is a second-difference operator. For a vector of length
$p$, the operator can be expressed as a banded matrix and applied via
\texttt{scipy.linalg.solve\_banded}. The compact bank includes one
Whittaker smoother with $\lambda=10^{3}$; the extended bank adds
$\lambda\in\{10,\,10^5\}$ to span a smoothing range. The bandedness is
exploited to keep $\Smat_b=\Amat_b\Smat$ at $\mathcal{O}(pq)$.

\subsection{Norris-Williams gap derivative}
The Norris-Williams operator~\cite{norris1976influence} is the
composition of a moving-average smoother and a finite-difference
derivative with a fixed gap. We expose two presets:
\begin{itemize}[leftmargin=2em]
  \item \texttt{NorrisWilliams(gap=5, smoothing=5, order=1)}: moderate
        smoothing, first derivative.
  \item \texttt{NorrisWilliams(gap=11, smoothing=5, order=1)}: longer
        gap, useful when peaks are broader.
\end{itemize}
The extended bank also contains an order-2 variant
\texttt{NorrisWilliams(gap=5, smoothing=5, order=2)}.

\subsection{Savitzky-Golay parameters}
Bank entries cover $(11,2,0)$, $(21,3,0)$, $(11,2,1)$, $(21,3,1)$, and
$(11,2,2)$ for $(\text{window},\text{polyorder},\text{deriv})$. These
parameter triples cover the practical range used in NIRS without
introducing redundant entries that the bank-selection mechanism would
ignore in any case.

\subsection{Detrend projections}
The detrend operator with degree $d$ projects each row onto the
orthogonal complement of the polynomial subspace
$\mathrm{span}\{1,t,\dots,t^d\}$ where $t$ is a normalised wavelength
axis. Bank entries are $d=1$ (linear detrend) and $d=2$ (quadratic
detrend). These are strict linear operators in our sense because the
projection matrix is fixed once $p$ is fixed.

%% =====================================================================
\section{Implementation notes}\label{app:impl}
%% =====================================================================

\paragraph{Banded vs.\ dense matrices.}
We never instantiate the dense $p\times p$ matrix when an operator
admits a cheap convolution implementation. Each operator class
overrides \verb|_apply_cov_impl| (apply $\Amat$ to columns of the
cross-covariance) and \verb|_adjoint_vec_impl| (apply $\Amat^\top$ to a
vector or matrix). Materialised matrices are computed only when the
operator is exercised by the materialised \NIPALS reference engine.

\paragraph{Backends.}
The numpy backend is the reference. The Torch backend mirrors the
numpy semantics with \verb|torch.float32| or \verb|torch.float64|
tensors and exposes the same engines (\verb|simpls_covariance|,
\verb|nipals_adjoint|, \verb|superblock_simpls|). The unit tests
assert numerical equivalence under tight tolerances.

\paragraph{Result schema.}
Every benchmark row carries the master \verb|RESULT_COLUMNS| schema
plus the AOM-specific diagnostic columns
\verb|selected_operator_sequence_json| and
\verb|selected_operator_scores_json|. This keeps result files
joinable across variants and, when the master CSV is available,
joinable to TabPFN reference rows.

\bibliographystyle{plain}
\bibliography{../manuscript/references}

\end{document}
